% =======================================================
% 6. RIEPILOGO
% =======================================================
\section{Riepilogo}
Nella seguente sezione viene riportato un riepilogo quantitativo dello stato dei requisiti individuati, suddivisi in base al loro grado di importanza.

\subsection{Requisiti obbligatori soddisfatti}
\begin{table}[ht]
	\centering
	\renewcommand{\arraystretch}{1.4}
	\begin{tabularx}{0.8\textwidth}{|l|X|}
		\hline
		\rowcolor{primarycolor!10}
		\textbf{Tipologia Requisito} & \textbf{Numero Requisiti Obbligatori} \\
		\hline
		Funzionali & [Num] \\
		\hline
		Di Qualità & [Num] \\
		\hline
		Di Vincolo & [Num] \\
		\hline
		Prestazionali & [Num] \\
		\hline
		\textbf{Totale Obbligatori} & \textbf{[TOT]} \\
		\hline
	\end{tabularx}
	\caption{Riepilogo requisiti obbligatori}
\end{table}

\subsection{Requisiti desiderabili soddisfatti}
\begin{table}[ht]
	\centering
	\renewcommand{\arraystretch}{1.4}
	\begin{tabularx}{0.8\textwidth}{|l|X|}
		\hline
		\rowcolor{primarycolor!10}
		\textbf{Tipologia Requisito} & \textbf{Numero Requisiti Desiderabili} \\
		\hline
		Funzionali & [Num] \\
		\hline
		Di Qualità & [Num] \\
		\hline
		Di Vincolo & [Num] \\
		\hline
		Prestazionali & [Num] \\
		\hline
		\textbf{Totale Desiderabili} & \textbf{[TOT]} \\
		\hline
	\end{tabularx}
	\caption{Riepilogo requisiti desiderabili}
\end{table}

\subsection{Requisiti opzionali soddisfatti}
\begin{table}[ht]
	\centering
	\renewcommand{\arraystretch}{1.4}
	\begin{tabularx}{0.8\textwidth}{|l|X|}
		\hline
		\rowcolor{primarycolor!10}
		\textbf{Tipologia Requisito} & \textbf{Numero Requisiti Opzionali} \\
		\hline
		Funzionali & [Num] \\
		\hline
		Di Qualità & [Num] \\
		\hline
		Di Vincolo & [Num] \\
		\hline
		Prestazionali & [Num] \\
		\hline
		\textbf{Totale Opzionali} & \textbf{[TOT]} \\
		\hline
	\end{tabularx}
	\caption{Riepilogo requisiti opzionali}
\end{table}